\documentclass[11pt,a4paper]{article}
\usepackage[T1]{fontenc}
\usepackage[utf8]{inputenc}
\usepackage{lmodern}
\usepackage[margin=27mm]{geometry}
\usepackage{amsmath,amssymb,amsthm,mathtools,booktabs,microtype,float}
\usepackage[hidelinks]{hyperref}
\usepackage{enumitem}
\hypersetup{pdftitle={An exact relaxation optimum in a one-parameter family of four quantum measurement bases},pdfauthor={Lluis Eriksson},pdfsubject={Quantum information; spectral gaps; randomized projective measurements},pdfkeywords={measurement channel, spectral gap, mutually unbiased bases, dimension six, exact algebra}}
\newtheorem{theorem}{Theorem}[section]
\newtheorem{proposition}[theorem]{Proposition}
\newtheorem{lemma}[theorem]{Lemma}
\newtheorem{corollary}[theorem]{Corollary}
\theoremstyle{remark}
\newtheorem{remark}[theorem]{Remark}
\DeclareMathOperator{\Tr}{Tr}
\DeclareMathOperator{\spec}{spec}
\DeclareMathOperator{\diag}{diag}
\newcommand{\id}{\mathrm{id}}
\newcommand{\C}{\mathbb C}
\newcommand{\E}{\mathcal E}
\newcommand{\T}{\mathcal T}
\newcommand{\HH}{\mathcal H_0}
\newcommand{\ones}{\mathbf 1}
\setlength{\emergencystretch}{2em}
\title{An exact relaxation optimum in a one-parameter family of four quantum measurement bases}
\author{Lluis Eriksson\\\small Independent researcher, Sweden}
\date{11 October 2026\\\small Working manuscript, version 1}
\begin{document}
\maketitle
\begin{abstract}
Four orthonormal bases define a quantum measurement channel by averaging their nonselective rank-one projective measurements. We determine the exact Hilbert--Schmidt relaxation rate of the Poisson semigroup generated by this channel minus the identity, for the one-parameter subfamily introduced by Raynal, L\"u and Englert in their study of approximately mutually unbiased bases in dimension six. A block decomposition gives the complete spectrum as explicit quadratic polynomials in the parameter. The largest relaxation rate within this subfamily, with equal measurement probabilities and unit total event rate, is exactly $73/100$, attained uniquely at $u=7/10$. This parameter does not maximize the mean squared chordal distance between the bases. An exact rational comparison exhibits a larger mean distance with a smaller relaxation rate. We also prove the $2/3$ ceiling for arbitrary mixtures of product-vector measurement bases on $\C^2\otimes\C^3$, which the construction exceeds. The optimality statement is restricted to the specified subfamily; no global four-basis optimality or new mutually unbiased basis construction is claimed.
\end{abstract}

\section{Introduction and scope}
Mutually unbiased bases are useful in quantum state determination and related measurement problems \cite{wootters,zhu}. In dimension six, Raynal, L\"u and Englert (RLE) constructed a two-parameter family of four bases and identified a one-parameter subfamily that contains their maximizer of the mean squared chordal distance \cite{rle}. Their family provides an explicit setting in which to compare geometric and dynamical objectives.

We use those existing bases, with their attribution retained throughout. The quantity optimized here is the slowest exponential decay rate of a continuous-time, randomized measurement process. The resulting exact spectrum yields a different optimum from the mean-distance objective. The general possibility that distinct objectives have distinct optimizers is not the contribution; the result is the explicit spectral formula, its sharp optimization on the RLE subfamily, and an exact comparison of the two objectives there.

Measurement-induced mixing is related in motivation to noncommuting block-measurement diffusion \cite{eriksson}, but the processes differ. The present generator is a fixed average of complete rank-one measurements minus the identity. It is not a small-angle product limit and contains no additional Hamiltonian. Measurement-channel spectra also occur in quantum verification \cite{zhu} and classical shadows \cite{ippoliti}. We use the standard operator and Gram-matrix framework rather than claim that framework as new.

Every optimality assertion below specifies the admissible family and fixes both the measurement probabilities and the total event rate. In particular, the theorem does not optimize over the full two-parameter RLE family, over all four bases in $\C^6$, or over physical implementation costs. The proof is analytic; exact symbolic calculations supplied with the manuscript provide reproducibility checks rather than a substitute for the proof. Independent scientific assessment has not been completed for this version.

\section{Measurement dynamics and its Gram matrix}
Let $B_j=(b_{j1},\ldots,b_{j6})$, $j=0,1,2,3$, be orthonormal bases of $\C^6$. Put $\Pi_{jk}=|b_{jk}\rangle\langle b_{jk}|$ and
\begin{equation}\label{eq:channel}
 \E_j(X)=\sum_{k=1}^6\Pi_{jk}X\Pi_{jk},\qquad
 \T=\frac14\sum_{j=0}^3\E_j,\qquad \mathcal L=\T-\id.
\end{equation}
At each event of a rate-one Poisson process, one chooses a basis uniformly, measures, and discards the outcome. The averaged evolution is $e^{t\mathcal L}$. Indeed,
\[
 e^{t(\T-\id)}=e^{-t}\sum_{n=0}^{\infty}\frac{t^n}{n!}\T^n,
\]
which is completely positive, trace preserving and unital. Each $\E_j$ is an orthogonal projection in the Hilbert--Schmidt inner product. Thus $\T$ is self-adjoint and positive semidefinite as a superoperator, with spectrum in $[0,1]$.

Let $\HH=\{X\in M_6(\C):\Tr X=0\}$. Define
\begin{equation}\label{eq:gap}
 \gamma=1-\lambda_{\max}(\T|_{\HH}).
\end{equation}
This definition gives $\gamma=0$ if a nonscalar stationary operator exists. For every state $\rho$,
\begin{equation}\label{eq:decay}
 \|e^{t\mathcal L}(\rho)-I_6/6\|_2
 \le e^{-\gamma t}\|\rho-I_6/6\|_2.
\end{equation}
The exponent is sharp: the largest eigenvalue on $\HH$ has a nonzero Hermitian eigenoperator, and a sufficiently small multiple can be added to $I_6/6$ to obtain a density matrix.

Write $v_{jk}=\Pi_{jk}-I_6/6$ and form the real symmetric $24\times24$ matrix
\begin{equation}\label{eq:gram}
 G_{(j,k),(\ell,r)}=\Tr(v_{jk}v_{\ell r})
       =|\langle b_{jk},b_{\ell r}\rangle|^2-\frac16.
\end{equation}
The frame operator associated with these vectors is $4\T|_{\HH}$. Its nonzero eigenvalues agree with those of $G$, including multiplicity. Therefore
\begin{equation}\label{eq:gramgap}
 \gamma=1-\frac14\lambda_{\max}(G).
\end{equation}
The distinction between the dimension-$24$ Gram matrix and the dimension-$35$ traceless operator space matters for zero eigenvalues.

\section{The RLE subfamily and its complete spectrum}
We give all bases explicitly to fix phases, ordering and the parameter convention. These matrices are the RLE construction \cite[Eqs.~(5)--(7)]{rle}, expressed in notation adapted to the channel.
Set
\[
 \omega=e^{2\pi i/3},\quad
 F=\begin{pmatrix}1&1\\1&-1\end{pmatrix},\quad
 Z=\diag(1,-1),\quad X=\diag(\bar x,x),\quad
 H=\begin{pmatrix}1&\omega t^2\\1&-\omega t^2\end{pmatrix},
\]
where $|x|=|t|=1$. A bar denotes entrywise complex conjugation. Define $U_0=I_6$ and
\begin{align}
 U_1&=\frac1{\sqrt6}\diag\!\left(X,i\bar\omega tZ\bar X^2,X\right)
 \begin{pmatrix}F&F&F\\ F&\omega F&\bar\omega F\\ H&\bar\omega H&\omega H\end{pmatrix},\label{eq:u1}\\
 U_2&=\frac1{\sqrt6}
 \begin{pmatrix}F&F&F\\ H&\omega H&\bar\omega H\\ H&\bar\omega H&\omega H\end{pmatrix},\label{eq:u2}\\
 U_3&=\frac1{\sqrt6}\diag\!\left(\bar X,\bar\omega\bar X,-itZX^2\right)
 \begin{pmatrix}F&F&F\\ H&\omega H&\bar\omega H\\ F&\bar\omega F&\omega F\end{pmatrix}.\label{eq:u3}
\end{align}
The columns of $U_j$ are $B_j$. Every displayed matrix is unitary; the unitarity of the central block matrices follows from $FF^*=HH^*=2I_2$ and $1+\omega+\bar\omega=0$.

For the subfamily studied here, choose
\begin{equation}\label{eq:parameter}
 \frac14\le u\le1,\qquad
 x=\sqrt{1-u}+i\sqrt u,\qquad
 q=\frac{1-2u}{\sqrt u},\qquad
 t=e^{-i\pi/3}\bigl(q+i\sqrt{1-q^2}\bigr).
\end{equation}
The condition $|q|\le1$ is equivalent to $1/4\le u\le1$. Equation~\eqref{eq:parameter} fixes a branch of the RLE relation $\cos(\theta_t+\pi/3)=\cos(2\theta_x)/\sin\theta_x$, with $u=\sin^2\theta_x$. Equivalent phase branches are not needed for the theorem.

\subsection{Overlap blocks}
For $j\ne\ell$, let $P_{j\ell}$ have entries $|(U_j^*U_\ell)_{kr}|^2$. The canonical basis is unbiased with each of the other three, so $P_{0j}=J_6/6$, where $J_n$ is the all-ones matrix.
The remaining probabilities depend on just three real numbers:
\begin{equation}\label{eq:abc}
 A=\frac43(1-u)^2,\qquad B=(1-2u)^2,\qquad
 C=\frac{(1-u)(4u-1)}3,\qquad A+B+4C=1.
\end{equation}
With
\[
 D=\begin{pmatrix}A&B\\B&A\end{pmatrix},\qquad
 \widetilde D=\begin{pmatrix}B&A\\A&B\end{pmatrix},\qquad O=CJ_2,
\]
direct multiplication gives
\begin{equation}\label{eq:blocks}
 P_{12}=\begin{pmatrix}D&O&O\\O&D&O\\O&O&D\end{pmatrix},\quad
 P_{23}=\begin{pmatrix}O&D&O\\O&O&D\\D&O&O\end{pmatrix},\quad
 P_{31}=\begin{pmatrix}O&O&\widetilde D\\\widetilde D&O&O\\O&\widetilde D&O\end{pmatrix}.
\end{equation}
For completeness, one may obtain the first entry in \eqref{eq:abc} from the RLE overlap amplitude $\alpha=4\sqrt{1-u}(1-\omega\bar t\sqrt u)$. Since $\omega\bar t=-q+i\sqrt{1-q^2}$,
\[
 \frac{|\alpha|^2}{36}
 =\frac49(1-u)\bigl((1+\sqrt u\,q)^2+u(1-q^2)\bigr)
 =\frac43(1-u)^2.
\]
The companion amplitude has squared modulus $36(1-2u)^2$; the other four probabilities in each row equal $C$. Their values follow either by the same multiplication or by row normalization. At $u=7/10$, equations~\eqref{eq:u1}--\eqref{eq:parameter} together with \eqref{eq:blocks} are also checked directly in the supplementary exact verifier.

\begin{theorem}[Complete Gram spectrum]\label{thm:spectrum}
Set
\begin{equation}\label{eq:sk}
 s=A-B=\frac{1+4u-8u^2}{3},\qquad
 k=1-6C=8u^2-10u+3.
\end{equation}
For the bases \eqref{eq:u1}--\eqref{eq:parameter}, the Gram spectrum, with multiplicities, is
\begin{equation}\label{eq:spectrum}
 \spec(G)=\{0^{(4)},1^{(5)},(1-2s)^{(3)},(1+s)^{(6)},
                       (1+2k)^{(2)},(1-k)^{(4)}\}.
\end{equation}
When eigenvalues coincide, their multiplicities add. In particular,
\begin{equation}\label{eq:rateformula}
 \gamma(u)=1-\frac{1+\max\{0,-2s,s,2k,-k\}}4.
\end{equation}
\end{theorem}
\begin{proof}
Each diagonal block of $G$ is $Q=I_6-J_6/6$. The canonical block has no off-diagonal couplings and contributes $0^{(1)}$ and $1^{(5)}$.

For the other three blocks, identify the six coordinates with $\C^3\otimes\C^2$ and split each two-coordinate block into $e_+=(1,1)/\sqrt2$ and $e_-=(1,-1)/\sqrt2$. Let
\[
 R=\begin{pmatrix}0&1&0\\0&0&1\\1&0&0\end{pmatrix},\qquad R^3=I_3.
\]
On the antisymmetric space $\C^3\otimes e_-$, $Q$ is the identity, and the blocks $P_{12},P_{23},P_{31}$ act respectively as $sI_3,sR,-sR^2$. Conjugating the three-basis block by $\diag(I_3,I_3,R^2)$ reduces it to
\[
 \begin{pmatrix}1&s&-s\\s&1&s\\-s&s&1\end{pmatrix}\otimes I_3.
\]
The scalar $3\times3$ factor has eigenvalues $1-2s$ once and $1+s$ twice, giving multiplicities three and six.

On the symmetric space, the constant vector in $\C^3$ gives a null vector for each basis block, contributing three more zeros. Its orthogonal complement has dimension two. There the three overlap blocks, after subtracting $J_6/6$, act as $kI_3,kR,kR^2$ restricted to that complement: this uses $A+B-2C=1-6C=k$. The same conjugation gives the scalar matrix
\[
 \begin{pmatrix}1&k&k\\k&1&k\\k&k&1\end{pmatrix}
\]
on each of two copies. Its eigenvalues are $1+2k$ once and $1-k$ twice. This accounts for all $24$ dimensions and proves \eqref{eq:spectrum}. Equation~\eqref{eq:rateformula} follows from \eqref{eq:gramgap}.
\end{proof}

\section{The exact optimum and the distance comparison}
\begin{theorem}[Sharp subfamily optimum]\label{thm:optimum}
For equal probabilities and unit total measurement rate, the largest relaxation rate within \eqref{eq:parameter} is $73/100$. It is attained uniquely at $u=7/10$.
\end{theorem}
\begin{proof}
Put $h(u)=\max\{0,-2s,s,2k,-k\}$. We prove $h(u)\ge2/25$ over the entire parameter interval. For $1/4\le u\le3/5$, the function $s$ is nonincreasing and
\[
 h(u)\ge s(u)\ge s(3/5)=13/75>2/25.
\]
For $3/5\le u\le7/10$,
\[
 -k-\frac2{25}=-8\left(u-\frac{11}{20}\right)
                         \left(u-\frac7{10}\right)\ge0,
\]
with equality in that interval only at $u=7/10$. For $7/10\le u\le1$,
\[
 -2s-\frac2{25}=\frac{16}{3}\left(u-\frac7{10}\right)
                                     \left(u+\frac15\right)\ge0,
\]
again with equality only at $7/10$. At that parameter, $s=-1/25$, $k=-2/25$ and $h=2/25$. Substitution into \eqref{eq:rateformula} gives $1-(27/25)/4=73/100$.
\end{proof}

\begin{corollary}[An explicit algebraic construction]\label{cor:construction}
Taking
\[
 x=\frac{\sqrt{30}+i\sqrt{70}}{10},\qquad
 q=-\frac{2\sqrt{70}}{35},\qquad
 t=\frac{1-i\sqrt3}{2}\bigl(q+i\sqrt{1-q^2}\bigr)
\]
in \eqref{eq:u1}--\eqref{eq:u3} gives four bases whose channel on $\HH$ has spectrum
\begin{equation}\label{eq:optspectrum}
 \left\{0^{(15)},\left(\frac{21}{100}\right)^{(2)},
 \left(\frac{24}{100}\right)^{(6)},
 \left(\frac14\right)^{(5)},
 \left(\frac{27}{100}\right)^{(7)}\right\}.
\end{equation}
\end{corollary}
\begin{proof}
Here $(A,B,C)=(3/25,4/25,9/50)$. Use Theorem~\ref{thm:spectrum} and the nonzero-spectrum correspondence after \eqref{eq:gram}. The Gram rank is $20$, so the traceless channel has nullity $35-20=15$.
\end{proof}

For comparison with the original geometric objective, define the mean squared chordal distance as in \cite{rle}:
\begin{equation}\label{eq:distance}
 \overline{D^2}=1-\frac1{30}\sum_{0\le j<\ell\le3}
                             \|P_{j\ell}-J_6/6\|_F^2.
\end{equation}
Each noncanonical pair has squared Frobenius defect
$6(A^2+B^2+4C^2)-1$, while the three canonical pairs have zero defect. Hence
\begin{equation}\label{eq:distancepoly}
 \overline{D^2}(u)=
 -\frac{448u^4-1024u^3+888u^2-352u+25}{30}.
\end{equation}
At the spectral optimum,
\begin{equation}\label{eq:distanceatopt}
 \overline{D^2}(7/10)=\frac{6239}{6250},\qquad
 \left.\frac{d\overline{D^2}}{du}\right|_{u=7/10}=-\frac{12}{625}\ne0.
\end{equation}
Thus this interior parameter cannot maximize the distance objective. Table~\ref{tab:comparison} gives a strict comparison using only rational values. The increase in distance in its first row over its second is exactly $38513/750000000$, whereas its relaxation rate is smaller by $29/20000$.
\begin{table}[H]
\centering
\begin{tabular}{@{}lll@{}}
\toprule
Parameter $u$ & Mean distance $\overline{D^2}$ & Relaxation rate $\gamma$\\
\midrule
$139/200$ & $748718513/750000000$ & $14571/20000$\\
$7/10$ & $6239/6250$ & $73/100$\\
\bottomrule
\end{tabular}
\caption{A strict reversal between the geometric and dynamical objectives. Both columns are to be maximized. The first row has larger distance; the second has faster relaxation.}\label{tab:comparison}
\end{table}
The distance-maximizing parameter identified in \cite{rle} has $u\simeq0.6945524$, the real root of $112u^3-192u^2+111u-22=0$. Substitution into our rate formula gives approximately $0.728425$. Neither the proof of Theorem~\ref{thm:optimum} nor the exact comparison in Table~\ref{tab:comparison} depends on this decimal evaluation.

\section{Product measurements and unrestricted four-basis bounds}
\begin{proposition}[Product-vector ceiling]\label{prop:product}
Let $\T$ be any probability-weighted mixture of complete orthonormal bases consisting of product vectors on $\C^2\otimes\C^3$. For the rate-one generator $\T-\id$, $\gamma\le2/3$. This includes orthonormal product bases that are not tensor products of two fixed local bases.
\end{proposition}
\begin{proof}
Consider the three-dimensional operator subspace
$V=\{A\otimes I_3:\Tr A=0\}$. The orthogonal projection of a rank-one product projector $\Pi=\rho_A\otimes\rho_B$ onto $V$ is
\[
 Q_V\Pi=(\rho_A-I_2/2)\otimes I_3/3,
 \qquad \|Q_V\Pi\|_2^2=\frac16.
\]
Every basis has six projectors, and its dephasing superoperator is $\sum_k|\Pi_k)(\Pi_k|$. Its compressed superoperator trace on $V$ is therefore one, and the same holds for every probability mixture. The largest eigenvalue of $\T|_{\HH}$ is at least this compressed trace divided by $\dim V$, namely $1/3$. Equation~\eqref{eq:gap} proves the claim.
\end{proof}
The ceiling is sharp. For example, average independently over the three Pauli measurement bases on the qubit and four mutually unbiased measurement bases on the qutrit. The local traceless eigenvalues are $1/3$ and $1/4$, and the correlation-sector eigenvalue is $1/12$, giving $\gamma=2/3$. These finite local constructions are standard \cite{wootters}. The rate $73/100$ is therefore $9.5\%$ larger than the product-vector ceiling, with the total event rate held fixed. This is a rate comparison, not a statement about gate counts, laboratory time per measurement or energetic cost.

Let $\gamma_4^\star(6)$ denote the supremum over all four bases with uniform probabilities in dimension six. The construction and the elementary projector bound give
\begin{equation}\label{eq:globalbounds}
 \frac{73}{100}\le\gamma_4^\star(6)\le\frac34.
\end{equation}
To see the upper bound, restrict each $\E_j$ to $\HH$, where it is a nonzero orthogonal projection. On a unit vector in the range of any one of them, the quadratic form of their sum is at least one. Thus $\lambda_{\max}(4\T|_{\HH})\ge1$.
Equality in this bound requires the four traceless projection ranges to be mutually orthogonal, equivalently four mutually unbiased bases; this standard connection is also present in the verification formulation of \cite{zhu}. We neither establish equality nor sharpen the unrestricted upper bound. In particular, $73/100$ is not asserted to be the global optimum.

\section{Reproducibility, attribution and limitations}
The accompanying \texttt{verify.py} constructs the four matrices at the algebraic optimum, checks unitarity exactly, forms their Gram matrix directly and checks its characteristic polynomial over the rationals. Independently, it builds the polynomial block matrix \eqref{eq:blocks}, verifies the full characteristic polynomial over $\mathbb Q[u]$, and checks the exact distance comparison and the polynomial inequalities in the optimization proof. It uses Python and SymPy; no floating-point tolerance is needed for these checks. The supplementary \texttt{verification.json} records the actual run. The supplement contains source and environment requirements, but the mathematical argument is self-contained in this manuscript.

The RLE bases, their parameter relation and their original distance objective are prior work. The projector-frame interpretation is also standard. A focused literature search did not identify the particular sharp subfamily relaxation optimum proved here, but that search is not an exhaustive priority certificate. The paper makes a narrowly stated contribution and leaves its novelty and significance open to independent assessment. No formal proof-assistant verification or independent human review is claimed.

\paragraph{AI assistance.}
This manuscript was prepared with substantial assistance from OpenAI Codex during a research session on 11 October 2026. The assistant performed literature searches, proposed the spectral optimization problem, derived the formulas, wrote the draft and verification code, and executed symbolic and numerical checks. The author commissioned the work and requested preparation for AIRR publication. Independent post-publication model assessments were requested as a subsequent step and have not been performed for this version. The AI assistance and computational checks should not be read as independent peer review.

\paragraph{Conflict of interest.}
The author operates AIRR.science. Hosting or release there is distinct from independent scientific assessment. This version is presented as a working manuscript, without an editorial acceptance or reviewer score.

\begin{thebibliography}{9}
\bibitem{wootters}
W.~K. Wootters and B.~D. Fields,
``Optimal state-determination by mutually unbiased measurements,''
\emph{Annals of Physics} \textbf{191}, 363--381 (1989).
\href{https://doi.org/10.1016/0003-4916(89)90322-9}{doi:10.1016/0003-4916(89)90322-9}.

\bibitem{rle}
P. Raynal, X. L\"u and B.-G. Englert,
``Mutually unbiased bases in six dimensions: The four most distant bases,''
\emph{Physical Review A} \textbf{83}, 062303 (2011).
\href{https://doi.org/10.1103/PhysRevA.83.062303}{doi:10.1103/PhysRevA.83.062303};
\href{https://arxiv.org/abs/1103.1025}{arXiv:1103.1025}.

\bibitem{zhu}
H. Zhu and M. Hayashi,
``Optimal verification and fidelity estimation of maximally entangled states,''
\emph{Physical Review A} \textbf{99}, 052346 (2019).
\href{https://doi.org/10.1103/PhysRevA.99.052346}{doi:10.1103/PhysRevA.99.052346};
\href{https://arxiv.org/abs/1901.09772}{arXiv:1901.09772}.

\bibitem{ippoliti}
M. Ippoliti,
``Classical shadows based on locally-entangled measurements,''
\emph{Quantum} \textbf{8}, 1293 (2024).
\href{https://doi.org/10.22331/q-2024-03-21-1293}{doi:10.22331/q-2024-03-21-1293};
\href{https://arxiv.org/abs/2305.10723}{arXiv:2305.10723}.

\bibitem{eriksson}
L. Eriksson,
``Noncommuting Block Measurements Generate Quantum Diffusion:
A Three-Obstruction Convergence Hierarchy, a Noncommutative Connectivity Gap, and Exact Algebraic Architectures,''
AIRR, ARR-2026-3A3YP2CEH79D4RV0 (2026).
\url{https://airr.science/papers/ARR-2026-3A3YP2CEH79D4RV0/}.
\end{thebibliography}
\end{document}
